% !PRINTED 450
\weberpage{450}
%<<CONT
\gesperrt{kann durch Substitutionen der Gruppe} $G_{168}$ \gesperrt{in jeden anderen transformirt werden.}

Setzen wir
\weq{3}{T = x_1 + x_2 + x_3,}
so ist $T = 0$ die Gleichung der Verbindungslinie der Punkte $\pi, \pi'$. Die Function $T$ bleibt durch $\chi$ ungeändert und wechselt durch $\omega$ ihr Vorzeichen. Sie geht aber durch die 28 Substitutionen $\Theta^s\tau^r$ in 28 \gesperrt{verschiedene} Functionen $T_{s,r}$ über, die wir mit Hülfe der Formeln des \S.\,115, (11), (14) leicht so darstellen können:
\weq{4}{
\begin{aligned}
T_{0,r} &= \varepsilon^r x_1 + \varepsilon^{2r} x_2 + \varepsilon^{4r} x_3 \\
h\,T_{1,r} &= \beta^2\varepsilon^{1+r} x_1 + \alpha^2\varepsilon^{2+2r} x_2 + \gamma^2\varepsilon^{4+4r} x_3 \\
-h\,T_{2,r} &= \alpha^2\varepsilon^{2+r} x_1 + \gamma^2\varepsilon^{4+2r} x_2 + \beta^2\varepsilon^{1+4r} x_3 \\
h\,T_{3,r} &= \gamma^2\varepsilon^{4+r} x_1 + \beta^2\varepsilon^{1+2r} x_2 + \alpha^2\varepsilon^{2+4r} x_3,
\end{aligned}
}
und diese stellen, gleich Null gesetzt, 28 verschiedene gerade Linien dar, die alle aus einer von ihnen durch Substitutionen der Gruppe ableitbar sind.

Auf der Linie $T$ liegen die drei Punkte mit den Coordinaten
\[
\begin{aligned}
\alpha\gamma &: \beta\gamma : \beta\alpha \\
\beta\alpha &: \alpha\gamma : \beta\gamma \\
\beta\gamma &: \beta\alpha : \alpha\gamma,
\end{aligned}
\]
die aus (10), \S.\,118 durch die Substitutionen $\tau^{-4}, \tau^{-4}\chi, \tau^{-4}\chi^2$ hervorgehen und also zu den achtzähligen Polen gehören.

6. \gesperrt{Auf jeder Linie $T$ liegen drei achtzählige und zwei dreizählige Pole.}

Es bleibt noch der dritte Pol $\pi_0$ der Substitution $\chi$ mit den Coordinaten
\[
x_1 = x_2 = x_3
\]
zu untersuchen, von dem wir schon nachgewiesen haben, dass er durch die sechs Substitutionen
\weq{5}{1, \chi, \chi^2; \omega, \omega\chi, \omega\chi^2}
ungeändert bleibt. Es fragt sich nun, ob dieser Pol nicht mehr als sechszählig ist.

Wenn es ausser (5) noch andere Substitutionen giebt, die den Pol $\pi_0$ unverändert lassen, so müssen diese nothwendig vom $2^{\mathrm{ten}}$ Grade sein, da $\pi_0$ weder ein vierzähliger noch ein siebenzähliger Pol ist, und weil ferner der Grad der zu $\pi_0$ gehörigen Gruppe nicht durch 9 theilbar sein kann.
